\documentclass[11pt,a4paper]{article}
\usepackage[utf8]{inputenc}
\usepackage[margin=1in]{geometry}
\usepackage{amsmath,amssymb}
\usepackage{booktabs}
\usepackage{tabularx}
\usepackage{hyperref}
\usepackage{cite}
\usepackage{microtype}
\usepackage{graphicx}

\title{\textbf{Classical Ge'ez NLP: A Curated 25,000-Sentence Benchmark Corpus and Multimodal Morpho-Syllabic Tokenizer for Under-Resourced Ethiopic Semitic}}

\author{
  \textbf{Meron Ghirmai} \\
  Independent Research in Computational Linguistics \& Semitic NLP \\
  \texttt{meronghirmai25@gmail.com} \\
  \url{https://github.com/MeronGhirmai/geez-nlp-corpus-tokenizer}
}

\date{September 2026}

\begin{document}

\maketitle

\begin{abstract}
Classical Ge'ez (ግዕዝ, ISO 639-2: \texttt{gez}) is an ancient South Semitic Afroasiatic language that serves as the liturgical, epigraphic, and historical literary bedrock of both Eritrea and Ethiopia. Despite its preservation across thousands of parchment codices dating from the 4th century CE to the 19th century CE, Ge'ez remains critically under-resourced in Natural Language Processing (NLP). Contemporary multilingual Large Language Models (LLMs) suffer from severe subword over-fragmentation, excessive fertility rates ($>3.8$ subwords/word), and syntactic blindness when processing Ethiopic scripts due to non-concatenative root-and-pattern morphology, proclitic agglutination, and the Fidel abugida syllabic structure.

In this work, we present the largest open-access, linguistically annotated Classical Ge'ez corpus to date, comprising \textbf{25,480 curated sentences ($\sim$354,210 tokens)} drawn from epigraphic stelae (Metera, Qohaito, Axum), royal and liturgical codices (Debre Bizen, Garima Gospels, Kebra Nagast, Book of Enoch), and classical grammatical chrestomathies (Dillmann, Chaîne, traditional monastic \textit{Sewasew}). Furthermore, we engineer a \textbf{Multimodal Morpho-Syllabic Tokenizer} tailored specifically to Ge'ez. Empirical evaluation demonstrates that our clitic-aware BPE tokenizer reduces fertility from \textbf{3.84 tokens/word (in XLM-RoBERTa)} to \textbf{1.31 tokens/word} ($65.8\%$ reduction), while preserving grammatical boundaries and triconsonantal root semantics. All datasets, morphological lexicons, and interactive workbench tools are publicly released under open-access licenses.
\end{abstract}

\textbf{Keywords:} Classical Ge'ez, Afroasiatic NLP, Low-Resource Languages, Fidel Abugida, Tokenization, Semitic Morphology, Eritrea, Ethiopia, Digital Humanities.

\section{Introduction}
Natural language processing for ancient and liturgical languages has experienced rapid advances for Greco-Roman, Classical Arabic, and Sanskrit traditions. However, the classical written heritage of the Horn of Africa---Classical Ge'ez---remains severely marginalized in modern computational linguistics.

Ge'ez is the venerated historical, theological, and scientific medium shared across Eritrea and Ethiopia. Its literary canon spans nearly two millennia, preserving unique works of world cultural heritage, such as the only complete intact manuscript tradition of the apocalyptic \textit{Book of Enoch} (መጽሐፈ ሄኖክ), the 5th-century \textit{Garima Gospels}, the 4th-century \textit{Metera} royal stele in Eritrea, and the monastic libraries of \textit{Debre Bizen} and \textit{Lake Tana}.

Modern language models (e.g., LLaMA, GPT-4, mBERT, XLM-RoBERTa) demonstrate severe degradation when processing Ge'ez text due to three fundamental issues:
\begin{enumerate}
    \item \textbf{Corpus Scarcity:} Available training sets contain virtually zero high-quality, normalized Classical Ge'ez text, frequently conflating Ge'ez with modern Amharic or Tigrinya.
    \item \textbf{Abugida Tokenizer Inefficiency:} BPE vocabularies trained predominantly on Latin-script corpora break down Ethiopic syllabographs into multiple raw UTF-8 byte sequences, yielding extreme fertility rates ($>3.8$ subwords per word) that waste context window capacity.
    \item \textbf{Agglutinative Proclitic Binding:} Classical Ge'ez routinely compounds conjunctions, prepositions, relative markers, and enclitic pronouns directly into the lexical noun or verb stem, blinding naïve whitespace tokenizers to verbal roots.
\end{enumerate}

To address these challenges, this paper presents:
\begin{itemize}
    \item \textbf{The Curated 25,480-Sentence Ge'ez Corpus:} A digitized, verified, and morphologically analyzed corpus of 25,480 sentences ($\sim$354,210 tokens) spanning the 4th through 19th centuries CE, aligned with English translations and Universal Dependencies annotations.
    \item \textbf{Bi-National Heritage Coverage:} Sourcing uniting both Eritrean witnesses (\textit{Metera}, \textit{Debre Bizen}, \textit{Debre Sina}) and Ethiopian codices (\textit{Garima}, \textit{Axum}, \textit{Gondar}).
    \item \textbf{A Custom Multimodal Ge'ez Tokenizer:} Integrating clitic decomposition, root-aware BPE, orthographic division, and Fidel abugida phonetic-vocalic order deconstruction.
    \item \textbf{An Open-Source Research Workbench:} An interactive web application featuring KWIC concordance, Fidel matrix inspection, community annotation tools, and direct Python/Hugging Face integrations.
\end{itemize}

\section{Linguistic Characteristics \& Computational Bottlenecks}
\subsection{The Fidel Abugida Writing System}
Unlike alphabetic systems (Latin, Greek) or pure abjads (Arabic, Hebrew), the Ge'ez script is an \textbf{abugida} (alphasyllabary). The foundational inventory consists of 33 base consonants across 7 vocalic orders (\textit{Hälgät}, \textit{Ka'eb}, \textit{Saləs}, \textit{Rab'e}, \textit{Haməs}, \textit{Sadəs}, \textit{Sab'e}), yielding 231 primary consonant-vowel combinations, complemented by labiovelar series and numerical glyphs:
\begin{equation}
\text{Glyph}(C, V) = f(C_i, O_j) \quad \text{where } i \in [1, 33], \, j \in [1, 7]
\end{equation}
In standard Unicode, each syllable is encoded as a single precomposed code point. Standard UTF-8 subword tokenizers treat each code point as 3 distinct bytes, completely missing the systematic phonological relationships between different orders of the same consonantal root.

\subsection{Non-Concatenative Root-and-Pattern Morphology}
Similar to other Semitic languages, Ge'ez verbs are derived from discontinuous triconsonantal roots (e.g., $k-t-b$ ``to write'', $n-g-s$ ``to reign'', $q-d-s$ ``to be holy'') mapped onto vocalic templates (e.g., Perfect \textit{kataba}, Imperfect \textit{yəkatteb}, Subjunctive \textit{yəktəb}).

\subsection{Proclitic Agglutination}
In classical manuscripts, prepositions and conjunctions are orthographically fused to their host words without spacing. For example:
\begin{equation}
\text{ወበደብረ} \rightarrow \text{ወ-} (\text{and}) + \text{በ-} (\text{at/in}) + \text{ደብር} (\text{mountain}) + \text{-a} (\text{construct state})
\end{equation}
A naive whitespace tokenizer interprets ``ወበደብረ'' as an atomic out-of-vocabulary word, rather than recognizing the core lemma ``ደብር'' prefixed by two high-frequency functional clitics.

\section{Corpus Construction \& Methodology}
\subsection{Corpus Composition \& Provenance}
The corpus was curated across four balanced genres totaling \textbf{25,480 sentences} and \textbf{354,210 tokens}:

\begin{table*}[t]
\centering
\small
\begin{tabular}{llrrc}
\toprule
\textbf{Genre} & \textbf{Primary Historical Sources} & \textbf{Sentences} & \textbf{Tokens} & \textbf{Share (\%)} \\
\midrule
Historical Manuscripts & Debre Bizen (Eritrea), Garima, Kebra Nagast, Enoch & 11,460 & 162,732 & 45.0\% \\
Educational Grammars   & Monastic \textit{Sewasew}, Dillmann, Weninger, Chaîne & 7,640 & 106,263 & 30.0\% \\
Monastic Submissions   & Debre Bizen, Debre Damo, Lake Tana, Qene Poetry     & 4,890 & 68,008 & 19.2\% \\
Epigraphic Stelae      & Metera (Eritrea), Qohaito, Axumite Royal Stelae    & 1,490 & 17,207 & 5.8\% \\
\midrule
\textbf{Total}         & — & \textbf{25,480} & \textbf{354,210} & \textbf{100.0\%} \\
\bottomrule
\end{tabular}
\caption{Corpus breakdown across classical genres, manuscripts, and epigraphy.}
\label{tab:corpus}
\end{table*}

\subsection{Data Cleaning \& Annotation Pipeline}
Raw transcriptions underwent a multi-stage normalization pipeline:
\begin{enumerate}
    \item \textbf{Unicode Canonicalization:} Replacing Latin punctuation with Ethiopic Unicode standards (\texttt{U+1361} for \texttt{፡}, \texttt{U+1362} for \texttt{።}).
    \item \textbf{Clitic Segmentation:} Disambiguating proclitic prefixes (\textit{wa-}, \textit{la-}, \textit{ba-}, \textit{za-}) from native initial root consonants.
    \item \textbf{Universal Dependencies Alignment:} Structuring sentences into CoNLL-U format with part-of-speech tags, morphological features, and bilingual English glosses.
\end{enumerate}

\section{Custom Multimodal Tokenizer Architecture}
To overcome the limitations of off-the-shelf tokenizers, we engineered a dedicated 4-tier tokenization pipeline:
\begin{enumerate}
    \item \textbf{Orthographic Punctuation Tokenizer:} Isolates traditional Ethiopic delimiters without stripping them, preserving clausal structure.
    \item \textbf{Morphological Clitic Segmenter:} Applies finite-state constraints to decouple proclitic prepositions/conjunctions and enclitic suffixes from stems.
    \item \textbf{Subword BPE Encoder:} Utilizes a calibrated 8,192-token vocabulary trained exclusively on normalized Classical Ge'ez corpora.
    \item \textbf{Fidel Abugida Decomposer:} Deconstructs each glyph into its base consonant and canonical vocalic order (1st to 7th order).
\end{enumerate}

\section{Quantitative Evaluation \& Benchmarks}
\subsection{Token Fertility Comparison}
Token fertility measures the average number of subwords produced per word:
\begin{equation}
\text{Fertility} = \frac{\text{Total Subword Tokens}}{\text{Total Words}}
\end{equation}

\begin{table}[h]
\centering
\small
\begin{tabular}{lrcc}
\toprule
\textbf{Tokenizer Model} & \textbf{Vocab Size} & \textbf{Fertility (Tokens/Word)} & \textbf{OOV Rate (\%)} \\
\midrule
mBERT (Multilingual BERT) & 119,547 & 3.62 & 2.14\% \\
XLM-RoBERTa               & 250,002 & 3.84 & 1.87\% \\
LLaMA 3 Tokenizer         & 128,256 & 4.12 & 0.00\% \\
\textbf{Our Base Ge'ez BPE} & \textbf{8,192} & \textbf{1.48} & \textbf{0.00\%} \\
\textbf{Our Clitic-Aware BPE} & \textbf{8,192} & \textbf{1.31} & \textbf{0.00\%} \\
\bottomrule
\end{tabular}
\caption{Token fertility and OOV rates on the Classical Ge'ez test benchmark.}
\label{tab:fertility}
\end{table}

As shown in Table~\ref{tab:fertility}, standard multilingual tokenizers generate nearly 4 subwords per word. Our clitic-aware BPE achieves \textbf{1.31 tokens/word}---yielding a \textbf{65.8\% reduction in token sequence length}.

\section{Open-Source Availability \& Reproducibility}
All resources developed in this project are publicly available:
\begin{itemize}
    \item \textbf{GitHub Repository:} \url{https://github.com/MeronGhirmai/geez-nlp-corpus-tokenizer}
    \item \textbf{Hugging Face Dataset:} \texttt{meronghirmai/classical-geez-corpus}
    \item \textbf{Hugging Face Tokenizer:} \texttt{meronghirmai/geez-bpe-tokenizer}
\end{itemize}

\section{Conclusion}
This research provides the first production-grade, open-source NLP corpus and specialized tokenizer for Classical Ge'ez. By providing 25,480 annotated sentences and reducing subword fertility by 65.8\%, this benchmark resolves fundamental computational hurdles for under-resourced Afroasiatic languages of the Horn of Africa.

\bibliographystyle{plain}
\begin{thebibliography}{9}
\bibitem{dillmann1899} A.~Dillmann, \textit{Grammatik der \"athiopischen Sprache}, T.O. Weigel, Leipzig, 1899.
\bibitem{chaine1907} M.~Cha\^ine, \textit{Grammaire \'Ethiopienne}, Imprimerie Catholique, Beyrouth, 1907.
\bibitem{leslau1987} W.~Leslau, \textit{Comparative Dictionary of Ge'ez (Classical Ethiopic)}, Harrassowitz, 1987.
\bibitem{sennrich2016} R.~Sennrich, B.~Haddow, and A.~Birch, ``Neural Machine Translation of Rare Words with Subword Units,'' \textit{ACL}, 2016.
\bibitem{kudo2018} T.~Kudo and J.~Richardson, ``SentencePiece: A Simple and Language Independent Subword Alternative for Neural Network Text Processing,'' \textit{EMNLP}, 2018.
\bibitem{ghirmai2026} M.~Ghirmai, ``Classical Ge'ez NLP Corpus and Multimodal Tokenizer for Under-Resourced Horn of Africa Languages,'' \textit{GitHub \& Hugging Face}, 2026.
\end{thebibliography}

\end{document}
